\section{Completed forms and their local parameters}\label{sec:completed}

In this section \(F,S,\chi,\psi\) have the meaning fixed in
Section~\ref{sec:preliminaries}. We construct completed forms and
identify their integral local-block action with the restriction of
the global Galois action. This comparison gives finiteness over the
local pseudodeformation base and describes the Banach coefficient
fibers needed for classicality. We then prove a lower bound on every
Hecke component and exact auxiliary-level identities preserving the
old deformation action; these are the inputs to patching. All local
groups at \(2\)
are \(\GL_2(\mathbb Q_2)\), since \(2\) splits completely in \(F\).

\subsection{Integral forms and classical density}

Fix an open tame level \(U^2\), hyperspecial outside \(S\).
We may also prescribe a finite-order transformation character
at the tame level, compatible with the central character.
Write \(\widetilde S_\psi(U^2,\OO)\) for the continuous
\(\OO\)-valued functions on
\[
 D_F^\times\backslash(D_F\otimes_F\mathbb A_F^\infty)^\times
\]
with the prescribed right tame transformation and central
transformation \(f(gz)=\psi(z)f(g)\). Its scalar extension to
\(E\) is a unitary Banach representation of
\[
 G_2=\prod_{v\mid2}\GL_2(F_v).
\]
The torsion version has values in \(E/\OO\) and is smooth.
The compact dual of the integral space is denoted by \(M\);
equivalently, it is the Pontryagin dual of the torsion space.
We also use this notation for a residual Hecke summand.

\begin{proposition}\label{prop:completed-forms}
For the tame level \(U^2\), central character \(\psi=\chi\eps\),
and full compact forms module \(M\) just defined, the following
properties hold. In the second assertion weights and local parameters
use the normalization of Section~\ref{sec:preliminaries}.
\begin{enumerate}
\item For a sufficiently small product \(K\) of uniform open
subgroups of \(\SL_2(F_v)\), the full compact module \(M\)
is finite free over \(\OO[[K]]\). The integral and torsion
function spaces are admissible and commute with reduction
of coefficients.
\item Locally algebraic vectors are classical definite
quaternionic forms of the corresponding algebraic weight.
Their regular noncharacter eigenpackets transfer to Hilbert
cuspidal eigenforms. They have Galois representations of
determinant \(\chi\), with the local compatibility described
below.
\item The regular classical vectors are dense in the Banach
space. They remain dense if, at each \(v\mid2\), one fixes a
smooth principal-series type with two distinct inertial
characters and varies the algebraic factor with the
prescribed central degree.
\item Every residual local factor \(\mathbb T\) of the
completed Hecke algebra away from \(S\) is a complete quotient
of \(R_F^{\mathrm{ps},\chi}\). It is reduced and
\(\OO\)-torsion-free, and its spectrum is the closure of its
classical coefficient points.
\end{enumerate}
\end{proposition}
\begin{proof}
Finiteness of quaternion ideal classes gives finitely many double
cosets at compact open level \cite[Main Theorem~17.7.1]{Voight2021};
passing to a smaller compact level has finite fibers. Stabilizers
modulo the center are finite as well. Their reduced norms have even
valuation at every finite place. The group
\[
 \{a\in F^\times:v(a)\in2\mathbb Z\text{ for every finite }v\}
       /(F^\times)^2
\]
is finite: the class of the half-valuation ideal maps it to
\(\operatorname{Cl}(F)[2]\), with kernel given by unit squareclasses,
and both groups are finite \cite[Theorems~4.4 and~5.1]{MilneANT}.
The kernel of reduced norm modulo squares has, after scalar
rescaling, norm-one representatives in a fixed compact subgroup;
these are finite by total definiteness
\cite[Lemma~17.7.13]{Voight2021}. Shrink \(K\) until
it meets none of these finite projective stabilizers
nontrivially. Choosing local sections for the central and
tame transformations identifies the function space with
continuous functions on a finite disjoint union of copies
of \(K\). Its compact dual is a finite direct sum of
\(\OO[[K]]\). This also proves admissibility and exact
coefficient reduction. Compatibility of the transformation
characters on their intersections is part of the definition;
the central character is totally even, since both \(\chi\)
and \(\eps\) are odd.

The identification of locally algebraic vectors is the
definite double-coset description of algebraic automorphic
forms. Jacquet--Langlands transfers the regular noncharacter
packets to Hilbert cusp forms
\cite[Theorems~14.4 and 16.1]{JacquetLanglands1970}; their
Galois representations, weights and de Rhamness are given by
\cite[Theorem~2.1.1]{BLGGT2014}. At the dyadic places we use the potentially
semistable comparison, including ramified types:
\cite[Corollary~1.2]{BLGGTLocalGlobal2014}. For the principal
types used here, the two distinct inertial characters
force monodromy to be zero. Thus the local representation
is potentially crystalline and crystabelian, with the
weights specified by the algebraic factor.

We spell out the density assertion with fixed type.
Let \(K_0=\prod_{v\mid2}\GL_2(\OO_{F_v})\).
Matrix coefficients of algebraic representations of
central degree zero form a subalgebra of the continuous
functions on \(K_0/Z(K_0)\). They contain constants and
separate points: the adjoint representation is faithful
on the projective group. Nonarchimedean polynomial
approximation, equivalently Stone--Weierstrass for this
compact totally disconnected space, makes this algebra
dense.

Choose a nonzero smooth type \(\tau\) and an algebraic
representation \(V\) with the needed central degree.
The matrix coefficients \(a_1,\ldots,a_N\) of
\(\tau\otimes V\) have no common zero, since every
representation matrix is invertible. For a continuous
section \(f\) with the prescribed central transformation,
choose a finite clopen partition subordinate to sets on
which some \(a_i\) is bounded away from zero. The functions
\(f/a_i\) on these clopens, extended by zero, are continuous
of central degree zero. Approximate them by the algebra
just described and multiply by \(a_i\).
Every resulting product is a matrix coefficient of
\(\tau\otimes V\otimes A\), where \(A\) is algebraic of
central degree zero. Complete reducibility in characteristic
zero decomposes \(V\otimes A\) into algebraic irreducibles
of exactly the same central degree. The smooth factor
\(\tau\) is unchanged. Averaging over the finitely many
projective stabilizers gives density in the double-coset
function space; its fixed denominators cause no problem
over \(E\).

At every dyadic factor we can choose the two inertial
characters to be distinct with the required product,
using characters of sufficiently large \(2\)-power order.
The associated compact principal-series type forces the
smooth local representation to be principal series of
that type. Here we use Henniart's compact types
\cite[Appendix~A, Sections~A.1.1--A.1.4 and A.2.3--A.2.4]{BreuilMezard2002},
whose scope includes residue cardinality two. In particular it excludes norm characters.
The algebraic central degree is common at all embeddings,
as required for classical weights. This proves the
assertion about a dense span of regular cuspidal packets.

\emph{The integral Hecke algebra and its residual factors.}
Form the Hecke algebra as the closure
of the integral Hecke operators in the endomorphisms of
the finite Iwasawa module. Reduction modulo its maximal
augmentation ideal gives a finite-dimensional commuting
algebra. Its topologically nilpotent kernel permits
lifting idempotents, giving finitely many residual factors.

\emph{Interpolation as an identity of integral operators.}
We next construct a continuous Galois trace in this integral Hecke
algebra. To pass from classical traces to actual operators, we first
specify finite tests for the operator topology. Work with
the full compact module
\(M\simeq\Lambda_K^r\), where
\(\Lambda_K=\OO[[K]]\) and
\(\mathfrak a=(\varpi,\ker(\Lambda_K\to\OO))\).
The topology on \(\End_{\Lambda_K}(M)\) is its
\(\mathfrak a\)-adic matrix topology. An operator \(A\)
is zero at depth \(j\) precisely when
\(A(M)\subseteq\mathfrak a^jM\), or equivalently
when its dual kills the finite \(\OO\)-module
\[
 Q_j=\Hom_\OO(M/\mathfrak a^jM,E/\OO).
\]
Choose finitely many generators of \(Q_j\). Exact
coefficient reduction in the function model represents
them by \(\varpi^{-e_i}f_i\) modulo integral functions,
with each \(f_i\) integral. Thus the finitely many tests
\(A^\vee f_i\in\varpi^{e_i}\widetilde S_\psi(U^2,\OO)\)
detect depth \(j\).

Approximate these \(f_i\), to the required precisions,
by finite sums of classical eigenvectors. Clear the
denominators of this finite collection once. Continuity
of the finitely many classical Galois traces then makes
the corresponding Hecke operators agree to those
precisions on the approximations whenever their Frobenius conjugacy
classes are sufficiently close. Integral Hecke operators
and their differences have norm at most one, so the
approximation errors are controlled uniformly. Hence
a Chebotarev sequence approaching any Galois conjugacy
class is Cauchy at every depth \(j\). Completeness gives
a unique integral limit, independent of the sequence;
the same finite tests prove continuity of the resulting
trace function. Restriction to the integral Hecke-idempotent
summands gives the assertion for each residual factor,
with its inherited operator topology.

\emph{The pseudodeformation quotient and its support.}
The determinant is the fixed character \(\chi\).
The polynomial identities for a two-dimensional
determinant hold on all classical packets and hence
on their dense span. Boundedness and torsion-freeness
make these identities integral. The universal property
gives a continuous map from
\(R_F^{\mathrm{ps},\chi}\) onto the completed trace
algebra: the Frobenius traces generate its dense
subalgebra, and the source is compact.

An operator vanishing at every classical eigenpacket
vanishes on the dense classical span, hence is zero.
Classical Hecke actions are semisimple. A nilpotent
operator therefore vanishes there, proving reducedness;
the action on integral functions proves
\(\OO\)-torsion-freeness. The same test proves the
asserted closure of the classical coefficient points.
\end{proof}

\subsection{Blocks, integral finiteness, and Banach fibers}

Fix a residual factor \(\mathbb T\). Extend \(k\) so that
the residual local parameters and the simple objects of
the corresponding local blocks are absolutely split.
For each \(v\mid2\), choose one representative from each simple
isomorphism class in the smooth block with central character
\(\psi_v\). Let \(P_v\) be the projective envelope of the direct
sum of their duals, taken in the compact dual category, and put
\(E_v=\End(P_v)\). This is the standard block generator of
\cite[Section~4.1]{PaskunasTung2021}.
The block parameter is the semisimplification of
\(\bar r_v\otimes\bar\eps\).
Denote the corresponding local pseudodeformation ring
by \(R_v^{\mathrm{ps}}\); its determinant is
\(\psi_v\eps=\chi_v\eps^2\).
Set
\[
 R_{\mathrm{loc}}^{\mathrm{ps}}
   =\widehat\bigotimes_{v\mid2,\OO} R_v^{\mathrm{ps}},
 \qquad
 E_{\mathrm{loc}}
   =\widehat\bigotimes_{v\mid2,\OO} E_v,
 \qquad
 P=\widehat\bigotimes_{v\mid2,\OO}P_v.
\]
Untwisting by \(\eps\) identifies this local base with the
fixed-\(\chi_v\) local parameters of the global problem.

\begin{proposition}\label{prop:block-comparison}
The residual Hecke summand belongs to the product block
just specified. The exact projective-generator functor gives
\[
 L=\Hom_{\mathrm{cts},G_2}(P,M),
\]
taken in the compact dual category of the specified product block,
with the right \(E_{\mathrm{loc}}\)-action by precomposition. It is
faithful over \(\mathbb T\), finite over
\(R_{\mathrm{loc}}^{\mathrm{ps}}\), and finite over
\(\mathbb T\). The local pseudodeformation action on \(L\)
is the restriction of the global Galois action, with the
stated cyclotomic twist. In particular,
\(\mathbb T\) is finite over \(R_{\mathrm{loc}}^{\mathrm{ps}}\).

At a coefficient point \(x:\mathbb T\to E'\), after finite
extension of \(E'\), its Banach Hecke eigenspace contains
\[
 \widehat\bigotimes_{v\mid2}\Pi_v,
\]
where \(\Pi_v\) is an irreducible unitary Banach
representation with local parameter
\((r_x|_{G_{F_v}}\otimes\eps)^{\mathrm{ss}}\), of determinant
\(\chi_v\eps^2\), in the block normalization of
Section~\ref{sec:preliminaries}. If every local parameter is absolutely
irreducible, de Rham, and regular, the eigenspace has a
nonzero locally algebraic vector. If the untwisted
parameter at \(v\) is
\(\gamma_{1,v}\oplus\gamma_{2,v}\), with ratio outside
\(\{1,\eps,\eps^{-1}\}\), the possibilities are the
continuous unnormalized upper-Borel inductions
\[
 \Ind_B^{\GL_2(\mathbb Q_2)}
       (\gamma_{i,v}\eps\otimes\gamma_{j,v}),
 \qquad \{i,j\}=\{1,2\}.
\]
\end{proposition}
\begin{proof}
\emph{The one-factor input.}
We first specify the one-factor results.
The integral center map and finiteness theorem of
Pa\v{s}k\=unas--Tung apply to every block, including the
scalar residual block at \(2\)
\cite[Theorems~1.1--1.4]{PaskunasTung2021}.
In particular \(E_v\) is finite over \(R_v^{\mathrm{ps}}\);
an integral-surjectivity claim about the center is
unnecessary. More precisely, integral finiteness is Theorem~1.3;
the center comparison in Theorem~1.4 can have integral cokernel killed
by~\(2\). Their compact projective-generator
equivalence is exact. Their generic reducible and
absolutely irreducible characteristic-zero fibers are
described in \cite[Corollary~6.10 and
Proposition~6.11]{PaskunasTung2021}.
The all-prime correspondence and locally algebraic
comparison are supplied by
\cite[Theorems~1.1 and 1.3]{CDP2014} and
\cite[Theorem~0.20]{Colmez2010}.
Thus the old residual-scalar restriction at \(2\) does not
enter the present application.

\emph{The product category.}
These statements extend to the required product category
by the argument of \cite[Lemmas~B.7--B.8]{GeeNewton2022},
adapted from the projective group to a fixed central
character using the all-prime inputs just stated.
Here are the details. Given a smooth admissible vector, fix compact
opens in all but one factor. Its cyclic span under the
remaining factor is finitely generated and admissible,
so is finite length by
\cite[Lemma~2.3.6 and Theorem~2.3.8]{Emerton2010OrdinaryI}.
The commuting actions give a surjection from the tensor
product of the one-factor cyclic spans onto the full
cyclic span. Absolutely split simple factors tensor to
simple product objects; filtering proves finite length.
Hence the smooth product representation is locally finite.

Apply the one-factor compact equivalences successively.
On dual finite-length objects this can be checked on simple
tensors and extended by exactness. Taking inverse limits
gives the equivalence for the compact locally finite
category, with endomorphism ring \(E_{\mathrm{loc}}\).
Admissibility gives finite multiplicities in the cosocle:
a homomorphism from a simple is detected by a nonzero
compact-fixed generating vector. There are only finitely
many simples in the block. Compact Nakayama therefore
makes \(\Hom(P,M)\) finite over \(E_{\mathrm{loc}}\).

\emph{Identifying the actual local and global actions.}
It remains to identify the block and its central action
with the global Galois restriction. Test one factor at a
time on the dense classical principal types of
Proposition~\ref{prop:completed-forms}. Local-global
compatibility gives a regular crystabelian Galois
parameter with distinct inertial characters.
For an irreducible parameter, the universal unitary
completion of its algebraic-smooth principal series is
the corresponding Banach representation by
\cite[Theorem~4.3.1 and Corollaries~5.3.1--5.3.4]{BergerBreuil2010}
and the locally algebraic comparison in \cite{CDP2014}.
Distinct inertia excludes both equal-root and special
degeneracies.

For an ordinary parameter with cyclotomic exponents
\(a>b\), the algebraic factor is
\[
 \operatorname{Sym}^{a-b-1} E^2\otimes\det^{b+1}.
\]
The ordinary completion theorem
\cite[Proposition~2.2.1]{BreuilEmerton2010} identifies the
relevant locally algebraic induction with a dense
subspace of the continuous unitary induction in the order
\[
 (\eta_l\eps,\eta_h),
\]
where \(\eta_h,\eta_l\) have exponents \(a,b\).
The following character calculation fixes the inducing order on the
dense ordinary test types; that order will matter again in
Section~\ref{sec:classicality}. For the precise normalization, put \(n=a-b-1\geq0\) and
\(\delta=(\eta_l\eps)\eta_h^{-1}\). Twisting by the full
unitary character \(\eta_h^{-1}\circ\det\) reduces this
induction to \(\Ind_B^{\GL_2(\QQ_2)}(\delta\otimes1)^{C^0}\).
With \(x\) the identity character and \(\eps=x|\cdot|\),
the source characters in the cited proposition, with \(k=n+2\), are
\[
 \chi_1=|\cdot|x^{-n},\qquad
 \chi_2=\delta|\cdot|^{-1}x^n.
\]
Since \(\delta\) is unitary of algebraic exponent \(-n\),
both \(\chi_2\) and \((\chi_2/\chi_1)x^{-n}\) are smooth, and
\[
 v_2(\chi_1(2))=-n-1,\qquad
 v_2(\chi_2(2))=n+1.
\]
The two characters in its completion formula are
\(\chi_1|\cdot|^{-1}x^n=1\) and
\(\chi_2|\cdot|x^{-n}=\delta\), in the reversed inducing
order \((\delta,1)\). Twisting back by
\(\eta_h\circ\det\) gives the displayed order and changes
the algebraic factor \(\det^{-n}\operatorname{Sym}^n\) into
\(\det^{b+1}\operatorname{Sym}^{a-b-1}\).
On each lower-coordinate disk its vectors are precisely
bounded-degree polynomial sections. The theorem concerns this
locally algebraic subspace, not the whole locally analytic induction.
When \(n=0\), the excluded equality
\(\chi_1=\chi_2|\cdot|^2\) is exactly \(\delta=1\),
which is ruled out by the distinct finite inertial characters.
The ordinary and nonordinary statements together supply
every test type used in the density argument; see also
\cite[Propositions~5.3--5.6]{PaskunasTung2021}.

A map from such a classical constituent into completed
forms extends continuously from its universal unitary
completion. After scaling lattices, it is a bounded map
in the compact block category. Naturality of the center
then identifies the central trace action on its image
with the Galois parameter given by classical
local-global compatibility. The normalizations give
\(\rho\otimes\eps\), as stipulated above.
The dense classical span proves equality of the two
actions on the whole Banach space. Torsion-freeness of
the integral lattice gives integral equality and shows
that every other block projection vanishes.

\emph{Integral finiteness and coefficient fibers.}
It follows from integral Pa\v{s}k\=unas--Tung finiteness
that \(L\) is finite over
\(R_{\mathrm{loc}}^{\mathrm{ps}}\).
Exactness and the equivalence imply that the action of
\(\mathbb T\) on \(L\) is faithful. The same finite
generators generate \(L\) over \(\mathbb T\).
Moreover,
\[
 \mathbb T\lhook\joinrel\longrightarrow
 \End_{R_{\mathrm{loc}}^{\mathrm{ps}}}(L);
\]
the right side is finite over the noetherian local base.
Thus \(\mathbb T\) is finite over that base.

At \(x\), the quotient of \(L[1/2]\) by the Hecke maximal
ideal is nonzero and finite-dimensional. Select a simple
quotient after a finite coefficient extension. The
commuting images of the factor endomorphism rings generate
a finite-dimensional algebra. Their radicals kill the
simple quotient; their central idempotents select
semisimple factors; after splitting those factors, the
quotient is a tensor product of simple modules for matrix
algebras. The original integral actions preserve a
lattice. All full finite lattices in this quotient are
commensurable. Under the inverse compact equivalence,
the tensor lattice corresponds to the completed tensor
product of the individual Banach lattices. Dualizing
the quotient therefore embeds the asserted Banach
tensor product in the eigenspace.

For an absolutely irreducible regular de Rham parameter,
the local factor has a nonzero locally algebraic vector,
so their tensor does as well. For a generic reducible
parameter, substituting
\(\gamma_i\eps\) into the fiber formula of
\cite[Corollary~6.10]{PaskunasTung2021} gives exactly the
two displayed inducing orders.
\end{proof}

\subsection{A lower bound on every Hecke component}

We next obtain a dimension bound that will supply the
initial support size for patching. For a finitely
generated compact \(k[[K]]\)-module, its growth degree is
the degree of growth of its length modulo powers of
the augmentation ideal.

\begin{lemma}\label{lem:simple-growth}
Let \(k\) be the characteristic-two residue coefficient field of
Section~\ref{sec:preliminaries}. The compact dual of an absolutely
irreducible admissible smooth \(k\)-representation of
\(\GL_2(\mathbb Q_2)\) with central
character has growth degree at most one on a sufficiently
deep determinant-one compact subgroup. Finite-length
objects for \(d\) factors have growth degree at most \(d\).
\end{lemma}
\begin{proof}
Write \(K_n=1+2^nM_2(\mathbb Z_2)\), \(n\ge2\).
Characters have bounded invariant dimension. For a
principal series, the double cosets indexed by
\(\mathbb P^1(\mathbb Z/2^n\mathbb Z)\) give
\(O(2^n)\) invariants. The Steinberg case follows from
the constants--principal-series--Steinberg exact
sequence: its possible \(H^1(K_n,k)\) correction is
bounded, since \(K_n\) has four generators.

For a supersingular \(\pi\), let
\[
 I_1=\begin{pmatrix}1+2\mathbb Z_2&\mathbb Z_2\\
                    2\mathbb Z_2&1+2\mathbb Z_2\end{pmatrix},
 \qquad \Pi=\begin{pmatrix}0&1\\2&0\end{pmatrix}.
\]
Thus \(I_1\) is the standard pro-\(2\) Iwahori subgroup.
Let \(M_\pi\) be the span of the translates of \(\pi^{I_1}\)
under \(\left(\begin{smallmatrix}2^n&b\\0&1\end{smallmatrix}\right)\),
where \(n\ge0\) and \(b\in\mathbb Z_2\);
this space is Iwahori-stable
\cite[Definition~4.5 and Lemma~4.6]{Paskunas2010}.
Use the exact sequence of Iwahori representations
\[
 0\longrightarrow \pi^{I_1}\longrightarrow
 M_\pi\oplus\Pi M_\pi\longrightarrow\pi\longrightarrow0
\]
from \cite[Theorem~1.2, proved as Theorem~6.3]{Paskunas2010}.
That statement explicitly includes \(p=2\).
By \cite[Proposition~4.7 and Lemma~4.10]{Paskunas2010},
\(M_\pi\) is the sum of two injective envelopes for the
upper unipotent group \(\mathbb Z_2\). Each has
\(2^n\)-dimensional invariants under \(2^n\mathbb Z_2\).
Conjugation by \(\Pi\) shifts the congruence level by
at most one. Since \(\dim\pi^{I_1}=2\), the invariants
sequence gives
\[
 \dim_k\pi^{K_n}\le 2\cdot2^n+2\cdot2^{n+1}+8.
\]
The last term bounds \(H^1(K_n,\pi^{I_1})\).
The classification of Barthel--Livn\'e and Breuil
\cite{BarthelLivne1994,BarthelLivne1995}, in the form of
\cite[Corollary~1.2]{Breuil2003I}, shows that the cases above exhaust
the irreducible admissible smooth representations with central character.

Let \(I\) be the augmentation ideal of a fixed deep
uniform compact group. In characteristic two, the
operators \(g^{2^n}-1=(g-1)^{2^n}\) belong to
\(I^{2^n}\). Thus the compact quotient modulo
\(I^{2^n}\) is bounded in dimension by the dual
\(K_{n+c}\)-invariants, for a fixed shift \(c\).
Choosing \(2^n\) comparable to an arbitrary exponent
gives linear augmentation growth.
The central character is trivial on a sufficiently
deep scalar subgroup, and the product of that subgroup
with a deep determinant-one subgroup has finite index.
The same upper bound therefore holds on the latter.
Finite extensions preserve upper growth bounds, and
tensor products add them. This proves the product
assertion.
\end{proof}

\begin{proposition}\label{prop:hecke-dimension}
For every irreducible component \(C\) of
\(\Spec\mathbb T\),
\[
 \dim C\ge1+2d.
\]
\end{proposition}
\begin{proof}
Reducedness and prime avoidance give an element of
\(\mathbb T\) that vanishes on every component except
\(C\) and is nonzero on \(C\). Its image on \(M\) is
a nonzero submodule \(M'\) supported on \(C\).
The full \(M\) is finite free over a sufficiently deep
Iwasawa algebra \(A=\OO[[K]]\).
The residual summand and \(M'\) embed in this free module.

The maximal augmentation filtration, including
\(\varpi\), has polynomial associated graded ring of
dimension \(1+3d\), after taking \(K\) sufficiently
deep. A nonzero submodule of a finite free \(A\)-module
has that full growth degree: it contains a cyclic
copy of \(A\), and Artin--Rees compares its induced
filtration with its intrinsic filtration.
Consequently \(M'\) has growth degree \(1+3d\).

Let \(r=\dim C\), and choose \(r\) parameters of the
supporting Hecke quotient. Write \(J\) for their ideal.
Proposition~\ref{prop:block-comparison} and exactness of
the block functor show that \(M'/JM'\) is dual finite
length. Lemma~\ref{lem:simple-growth} bounds its growth
degree by \(d\).
Each Hecke parameter is topologically nilpotent; on the
finite quotient \(M'/\mathfrak m_A M'\), some fixed
power is zero. Hence there is \(c>0\) such that
\[
 J^{cn}M'\subseteq\mathfrak m_A^nM'
 \quad(n\ge1).
\]
Filter \(M'/\mathfrak m_A^nM'\) by \(J\).
The degree-\(j\) layer is a quotient of at most
\(\binom{j+r-1}{r-1}\) copies of
\[
 (M'/JM')/\mathfrak m_A^n(M'/JM').
\]
Only \(cn\) layers occur. Their total length is therefore
\(O(n^{r+d})\).
Comparing with the degree \(1+3d\) gives
\(1+3d\le r+d\), which is the claim.
\end{proof}

\subsection{Auxiliary level and solvable base change}

The dimension bound supplies the initial support size. Patching
also requires that removing auxiliary level recover the original
operators, not just modules of the same rank. The following exact
statements connect completed forms
to the finite stages of localized patching. An auxiliary
prime has residue cardinality \(q_v\equiv1\pmod{2^n}\).
Inside its Iwahori group, take the kernel of the
diagonal-ratio map to a cyclic quotient of residue units
of order \(2^n\). This subgroup contains the scalars.
For \(h\) auxiliary primes, the deck group is
\(\Delta=(\mathbb Z/2^n)^h\).

\begin{lemma}\label{lem:auxiliary-level}
Assume that \(\chi\) is unramified at the auxiliary
primes. Localize at a characteristic-two point of the
completed Hecke spectrum at which their Frobenii have
distinct eigenvalues. The point may be the generic
point of an enhanced curve, rather than the original
finite-field residual point. The following assertions
hold for the localized completed-forms block modules,
and continue to hold after the coefficient extensions,
completions, and smooth framings used in
Section~\ref{sec:propagation}.
\begin{enumerate}
\item The auxiliary-level module is flat over
\(\OO[\Delta]\). Diamond augmentation is the Iwahori
module; after localization at the indicated point it is the
direct sum of \(2^h\) copies of the old hyperspecial
module, with its Galois-ring action factoring through
the old unramified quotient.
\item If \(U_j\) is a diamond generator and
\(\tau_j\) is the corresponding tame inertia element,
then
\[
 U_j^2-T(\tau_j)U_j+1=0.
\]
Congruent tame transformation characters give
isomorphic residue modules with the same deformation
actions.
\item A quadratic character trivial above \(2\) and at
infinity, unramified away from the auxiliary primes,
and factoring through the auxiliary character
quotient, acts by twisting forms. It gives an actual
semilinear isomorphism of the block module with its
Galois twist, preserving the central character, the
level, and the characteristic-two completion point.
\end{enumerate}
\end{lemma}
\begin{proof}
The finite disjoint-union function model used in
Proposition~\ref{prop:completed-forms} is free under the
deck action, since the ratio quotient kills scalars
and the level was chosen to remove projective
stabilizers. Its compact dual is therefore flat over
\(\OO[\Delta]\), with augmentation equal to the
Iwahori-level module. Apply the exact projective
functor to the finite-presentation ideal tests for
flatness. This proves flatness of the block module
and commutation of reduction with that functor.

At one auxiliary prime, let
\(\alpha:M_{\mathrm{old}}^{\oplus2}\to M_{\mathrm{Iw}}\)
be the two degeneracy maps and
\(\beta:M_{\mathrm{Iw}}\to M_{\mathrm{old}}^{\oplus2}\)
their trace maps. Their composite \(C=\beta\alpha\)
has diagonal entries \(q_v+1\) and off-diagonal entries
given by the unramified Hecke operator, up to central
units. Write \(D=\det C\). On the Iwahori module,
\(D\) denotes the same scalar expression in the
interpolated Frobenius trace and central operators.
In particular this is one compatible scalar operator
on both levels, not a determinant defined only on
the old module. The standard degeneracy calculation
expresses it, up to a central unit, as
\[
 (q_v+1)^2-q_v\chi(\phi_v)^{-1}T(\phi_v)^2,
\]
where \(\phi_v\) is a Frobenius lift. At the chosen
characteristic-two point, its value is a unit times
the nonzero square of the Frobenius trace. Thus
\(D\) becomes invertible in this localization.

Before inverting \(D\), we have the integral identity
\[
 \alpha\operatorname{adj}(C)\beta
       =D\,\operatorname{id}_{M_{\mathrm{Iw}}}.
\]
The local alternatives here are the Iwahori component of Henniart's
description \cite[Appendix~A, Sections~A.1.3 and~A.2.3]{BreuilMezard2002}.
Indeed an Iwahori-spherical classical cuspidal local
representation is unramified principal series or
special. In the former case its two Iwahori vectors
are spanned by the degeneracies, and the identity
follows from the matrix adjugate identity. In the
latter case the old space is zero and the Frobenius
eigenvalue ratio is \(q_v^{\pm1}\), so the displayed
scalar expression for \(D\) is zero as well.
Classical density and torsion-freeness prove the
identity on integral completed forms and their
compact duals. Consequently
\(D^{-1}\operatorname{adj}(C)\beta\) is both a
left and a right inverse to \(\alpha\) after the
selected localization. Iterating gives \(2^h\) old
copies.

The maps are Hecke equivariant. Interpolation
identifies their Galois trace actions; at the
absolutely irreducible enhanced curve, the framed
representation action is consequently the actual
old unramified action. Thus the augmentation
factors through the old deformation quotient,
with its original action, rather than only having
the same multiplicity. The flatness and Iwahori
augmentation established before localization
remain valid at the selected point on both sides.
The exact block functor and the finite-ideal tests
used in Section~\ref{sec:propagation} preserve them
under the indicated flat coefficient extensions,
localizations, completions, and smooth framings.

On a classical point with a diamond eigencharacter,
local Langlands gives semisimple inertia characters
equal to that character and its inverse, because
\(\chi\) is unramified. Their sum is \(T(\tau_j)\).
The displayed monic quadratic identity follows at
every such point, and hence on the integral module
by density. Fixed tame character transformations
congruent modulo \(\varpi\) identify their residue
function spaces. The associated trace comparisons
show that this identification also preserves the
residue deformation actions.

For the last assertion, multiply a form by the
quadratic character composed with reduced norm.
The central character is unchanged because its value
on a scalar norm is a square. On the ratio-defined
auxiliary subgroup, the two residue diagonal entries
have the same image in the character quotient; their
product is therefore killed by a quadratic character.
Thus the level is preserved. Triviality at \(2\) and
infinity preserves the local actions there, and the
Galois parameter twists by the same character.
Its residual reduction is one, so it fixes the
characteristic-two point and its residual frames.
This gives the asserted semilinear intertwiner.

Finally, adding a coefficient variable preserves
flatness by the complete-module local criterion.
Localization, completion of finite noetherian
modules, and smooth framing extension preserve the
stated flatness and module identities. The
intertwiners are local and preserve maximal-ideal
filtrations. Hence every assertion survives the
operations used at the finite patching stages.
\end{proof}

\begin{lemma}\label{lem:forward-basechange}
Let \(F'/F\) be a solvable totally real Galois extension split
completely at \(2\). The restriction to \(F'\) of a pro-modular
trace locus over \(F\) is pro-modular over \(F'\), at one common
tame level for the whole restricted locus.
If all its bad inertial types are killed, this level
can be taken Iwahori at the bad non-dyadic primes
and hyperspecial elsewhere away from \(2\).
For a general solvable totally real extension split completely
at \(2\), the restricted locus is potentially pro-modular.
\end{lemma}
\begin{proof}
First suppose that \(F'/F\) is Galois. Transfer the dense classical
points by solvable base change, iterated
through prime-degree cyclic extensions
\cite[Chapter~III, Theorems~4.2 and~5.1]{ArthurClozel1989}; for the
original \(\GL_2\) construction see
\cite[Chapter~11, Lemma~11.3 and Proposition~11.4]{Langlands1980}. Apply
the unitary formulation after a norm twist when necessary, then undo
the twist. Use
Jacquet--Langlands \cite[Theorems~14.4 and 16.1]{JacquetLanglands1970}.
A regular
cuspidal packet stays cuspidal along each prime-degree cyclic
step: a loss of cuspidality would make it induced
from a character of a totally real quadratic field,
whose parallel algebraic exponents contradict
regularity. Local-global compatibility and the
uniform inertia control of
Lemma~\ref{lem:inertia-basechange} bound the
conductors at a fixed finite set of places.
The wild kernel is common and the rank is two, so
the possible tame monodromy contribution is
uniformly bounded. Thus a single tame level
contains every transferred packet.

Their trace closure contains the restricted locus
by density and continuity. Thus this locus is pro-modular over
the specified field \(F'\), at the common tame level just chosen.
After the inertial types are killed, only unramified principal or
special local factors remain at a bad
non-dyadic prime, both with Iwahori invariants.
The stated level follows.

For a general solvable extension, pass to its solvable Galois
closure. It is still totally real and completely split at \(2\).
The Galois case then proves potential pro-modularity of the
restricted locus. We use only this forward
base-change assertion here; classical descent
is applied later to an actual cuspidal packet.
\end{proof}
